\section{Explicit Response Formula}
\quad We write down the explicit response formulas corresponding to the steady-state linear response theory developed in the previous section. As an example, we consider the response of the equal-time covariance matrix $\mathbf{K}_0 \equiv \langle \delta \vec{x} \delta \vec{x}^\intercal \rangle$ with respect to the change in $\vec{p}$, and $\vec{p}$ can be $\vec{r}$, $\mathbf{W}$, and $\boldsymbol{\sigma}$. The observable becomes $B_{ij}(\vec{x}; \vec{p})=\delta x_i \delta x_j$. We also present the corresponding response functions for the different types of parameter perturbations. From the general formula (\ref{response_function2}), the response function with respect to $p_\alpha$ gives
\begin{equation}
    \chi_{ij\alpha} = \langle \partial_p (\delta x_i \delta x_j) \rangle_{ss, \vec{p}} + \langle \delta (\delta x_i \delta x_j) \delta \mathscr{F}_\alpha \rangle_{ss, \vec{p}}.
\end{equation}
The observable is written as $\delta x_i \delta x_j= [x_i - \langle x_i \rangle][x_j - \langle x_j \rangle]$. Then
\begin{equation}
    \left\langle \frac{\partial B_{ij}}{\partial p_\alpha} \right\rangle = -\frac{\partial \langle x_i \rangle}{\partial p_\alpha} \langle \delta x_j \rangle - \frac{\partial \left\langle x_j \right\rangle}{\partial p_\alpha} \langle \delta x_i \rangle.
\end{equation}
Since $\langle \delta x_i \rangle = 0$, the first term vanishes. The second term can be written as
\begin{eqnarray}
    \left\langle \delta (\delta x_i \delta x_j) \delta \mathscr{F}_\alpha \right\rangle_{ss, \vec{p}} &=& \left\langle (\delta x_i \delta x_j) \mathscr{F}_\alpha \right\rangle_{ss, \vec{p}} - \left\langle \delta x_i \delta x_j \right\rangle_{ss, \vec{p}} \left\langle \mathscr{F}_\alpha \right\rangle_{ss, \vec{p}}\nonumber\\
    &=& \langle (\delta x_i \delta x_j) \mathscr{F}_\alpha \rangle_{ss, \vec{p}} - (\mathbf{K}_0)_{ij} \langle \mathscr{F}_\alpha \rangle_{ss, \vec{p}}.
\end{eqnarray}
The explicit formulas of $\mathscr{F}_\alpha$ for three kinds of the network parameters are shown in Appendix~B.
\subsection{Response due to Change in the Intrinsic Node Parameter $r$}
\quad We first consider a perturbation of a single node perturbation $r_\alpha$, where $\alpha$ denotes the index of the perturbed node. One obtains the steady-state response of the equal-time covariance matrix
\begin{equation}
    \chi_{ij\alpha} = \langle \delta x_i \delta x_j \mathscr{F}_\alpha \rangle - (\mathbf{K}_0)_{ij} \langle \mathscr{F}_\alpha \rangle.
    \label{response_Ko_r}
\end{equation}
By using the dependence of $\mathbf{Q}$ on $r_\alpha$, the response function can be written as
\begin{eqnarray}
    \chi_{ij\alpha}^{(ss)} &=& \frac{\partial f'_\alpha}{\partial r_\alpha} \Big|_{\vec{X}}\left( \frac{\partial \mathbf{K}_0}{\partial Q_{\alpha\alpha}} \right)_{ij}-\frac{\partial f_\alpha}{\partial r_\alpha} \Big|_{\vec{X}} \sum_s Q^{-1}_{s\alpha} \Bigg[ f_s''\left( \frac{\partial \mathbf{K}_0}{\partial Q_{ss}} \right)_{ij}\nonumber\\
    &+& \sum_n W_{sn} \partial_1^2 h(X_s, X_n) \left( \frac{\partial \mathbf{K}_0}{\partial Q_{ss}} \right)_{ij}+\sum_\mu W_{\mu s} \partial_1 \partial_2 h(X_\mu, X_s) \left( \frac{\partial \mathbf{K}_0}{\partial Q_{\mu\mu}} \right)_{ij}\nonumber\\
    &+& \sum_\nu W_{s\nu} \partial_1\partial_2 h(X_s, X_\nu) \left( \frac{\partial \mathbf{K}_0}{\partial Q_{s\nu}} \right)_{ij} + \sum_\mu W_{\mu s} \partial^2_2 h(X_\mu, X_s) \left( \frac{\partial \mathbf{K}_0}{\partial Q_{\mu s}} \right)_{ij} \Bigg].
\end{eqnarray}
For the simple cases of $f_i(x_i) = r_i(a_i-x_i)$ with linear $h$, or Logistic $f_i(x_i) = r_i x_i (1-x_i)$ with synchronizing coupling and $h(X_i, X_j) = 0$, the above result simplifies significantly to 
\begin{equation}
    \chi_{ij\alpha}^{(ss)} = -\left(\frac{\partial \mathbf{K}_0}{\partial Q_{\alpha\alpha}} \right)_{ij}.
    \label{chiKodr_ss}
\end{equation}
This form shows that the response function of the covariance matrix is directly controlled by the diagonal element $Q_{\alpha\alpha}=-r_\alpha-\sum_j W_{\alpha j}$, namely by the local restoring parameter and the total incoming coupling associated with node $\alpha$. The derivative of $\mathbf{K}_0$ can be solved by differentiating the two sides of the Lyapunov equation (\ref{Lyapunov_eq}), which will be discussed later in the section. 

\subsection{Response due to Change in the Connection Weight $W$}
\quad We next consider the perturbation of a directed connection weight $W_{\alpha\beta}$. The corresponding response function of the covariance matrix is
\begin{equation}
    \chi_{ij(\alpha\beta)} = \langle \delta x_i \delta x_j \mathscr{F}_{\alpha\beta} \rangle - (\mathbf{K}_0)_{ij} \langle \mathscr{F}_{\alpha\beta} \rangle.
    \label{response_Ko_W}
\end{equation}
Again, by using the dependence of $\mathbf{Q}$ on $W_{\alpha\beta}$, the explicit formula is
\begin{eqnarray}
    \chi_{ij\alpha\beta}^{(ss)} &=& \partial_1 h(X_\alpha, X_\beta) \left( \frac{\partial \mathbf{K}_0}{\partial Q_{\alpha\alpha}} \right)_{ij} + \partial_2 h(X_\alpha, X_\beta) \left( \frac{\partial \mathbf{K}_0}{\partial Q_{\alpha\beta}} \right)_{ij}\nonumber\\
    &-& h(X_\alpha,X_\beta) \sum_s Q^{-1}_{s\alpha} \Bigg[ \Bigg( f''_s + \sum_n W_{sn} \partial^2_1 h(X_s, X_n) \Bigg) \left( \frac{\partial \mathbf{K}_0}{\partial Q_{ss}} \right)_{ij}\nonumber\\
    &+& \sum_\mu W_{\mu s} \partial_1\partial_2 h(X_\mu, X_s) \left( \frac{\partial \mathbf{K}_0}{\partial Q_{\mu\mu}} \right)_{ij} + \sum_\nu W_{s\nu} \partial_1\partial_2 h(X_s, X_\nu) \left( \frac{\partial \mathbf{K}_0}{\partial Q_{s\nu}} \right)_{ij}\nonumber\\
    &+& \sum_\mu W_{\mu s} \partial_2^2 h(X_\mu, X_s) \left( \frac{\partial \mathbf{K}_0}{\partial Q_{\mu s}} \right)_{ij} \Bigg].
\end{eqnarray}
For systems that can achieve a fully synchronized noise-free state with synchronizing coupling function $h(X_i, X_j)=0$, the formula reduces to
\begin{equation}
    \chi_{ij(\alpha\beta)} = \partial_1 h(X_\alpha,X_\beta) \left( \frac{\partial \mathbf{K}_0}{\partial Q_{\alpha\alpha}} \right)_{ij} + \partial_2 h(X_\alpha,X_\beta) \left( \frac{\partial \mathbf{K}_0}{\partial Q_{\alpha\beta}} \right)_{ij},
    \label{chiKodW_ss}
\end{equation}
which makes the role of the perturbed connection especially transparent. In this case, the response is determined by two contributions: the change in the local diagonal part $Q_{\alpha\alpha}$, and the change in the directed off-diagonal coupling $Q_{\alpha\beta}$. For an undirected network, since $W_{\alpha\beta}=W_{\beta\alpha}$, the total response is the sum of the two directed contributions associated with $W_{\alpha\beta}$ and $W_{\beta\alpha}$.

\subsection{Response due to Change in the Noise Strength $\sigma$}
\quad Finally, we consider the perturbation of a noise strength. In contrast to the previous two cases, a change in $\sigma_{\alpha\alpha}$ does not alter the deterministic force, and therefore does not shift the fixed point. The corresponding response function gives
\begin{equation}
    \chi_{ij(\alpha\alpha)} = \langle \delta x_i \delta x_j \mathscr{F}_{\alpha\alpha} \rangle - (\mathbf{K}_0)_{ij} \langle \mathscr{F}_{\alpha\alpha} \rangle = \left( \frac{\partial \mathbf{K}_0}{\partial \sigma_{\alpha\alpha}} \right)_{ij}.
\end{equation}
This result emphasizes that a perturbation of the noise strength affects the steady-state fluctuations directly, without changing the deterministic fixed point.

\subsection{Evaluation of $\partial \mathbf{K}_0/\partial Q_{\alpha\beta}$ and $\partial \mathbf{K}_0/\partial \sigma_{\alpha\alpha}$}
\quad The derivative of $\mathbf{K}_0$ with respect to $\mathbf{Q}$ or $\boldsymbol{\sigma}$ is determined from the steady-state Lyapunov equation due to the dependence of $\mathbf{K}_0(\mathbf{Q},\boldsymbol{\sigma})$. The derivatives of $\mathbf{K}_0$ with respect to the matrix elements of $\mathbf{Q}$ can be obtained by differentiating the Lyapunov equation
\begin{equation}
    \mathbf{Q} \frac{\partial \mathbf{K}_0}{\partial Q_{\alpha\beta}} + \frac{\partial \mathbf{K}_0}{\partial Q_{\alpha\beta}} \mathbf{Q}^\intercal = \mathbf{M}^{(\alpha,\beta)},\; (\mathbf{M}^{(\alpha,\beta)})_{ij} \equiv -(\mathbf{K}_0)_{i\beta} \delta_{j\alpha} - (\mathbf{K}_0)_{j\beta} \delta_{i\alpha}.
    \label{derivative_Ko}
\end{equation}
Also, the derivative of the inverse of the covariance matrix gives
\begin{equation}
    \mathbf{K}_0 \frac{\partial \mathbf{K}^{-1}_0}{\partial Q_{\alpha\beta}} \mathbf{K}_0 = - \frac{\partial \mathbf{K}_0}{\partial Q_{\alpha\beta}}\;\text{or}\;\mathbf{K}^{-1}_0 \frac{\partial \mathbf{K}_0}{\partial Q_{\alpha\beta}} \mathbf{K}^{-1}_0 = - \frac{\partial \mathbf{K}^{-1}_0}{\partial Q_{\alpha\beta}}.
    \label{inverse_partialKo}
\end{equation}
Higher-order derivatives can be obtained by the similar manner. For the derivative with respect to $\sigma_{\alpha\alpha}$, we have
\begin{equation}
    \mathbf{Q}\frac{\partial\mathbf{K}_0}{\partial\sigma_{\alpha\alpha}} + \frac{\partial\mathbf{K}_0}{\partial\sigma_{\alpha\alpha}}\mathbf{Q}^\intercal = -\mathbf{E}^{(\alpha\alpha)},\;\text{where}\;E^{(\alpha\alpha)}_{ij}=1\;\text{for}\;i=j=\alpha.
    \label{partialKo_sigma}
\end{equation}
Moreover, the above formulas are further simplified when the network satisfies the equilibrium condition $\mathbf{Q}\boldsymbol{\sigma}=\boldsymbol{\sigma}\mathbf{Q}^\intercal$. In that case, the covariance matrix satisfies $\mathbf{K}_0 = -\frac{1}{2} \mathbf{Q}^{-1} \boldsymbol{\sigma}$, and therefore
\begin{equation}
    \left( \frac{\partial \mathbf{K}_0}{\partial Q_{\alpha\beta}} \right)_{ij} = -Q^{-1}_{i\alpha} (\mathbf{K}_0)_{j\beta} = \frac{1}{2} Q^{-1}_{i\alpha} (\mathbf{Q}^{-1}\boldsymbol{\sigma})_{j\beta}.
\end{equation}
If the nodes share the same noise strength $\boldsymbol{\sigma}=\sigma \mathbf{I}$, the formula is further simplified
\begin{equation}
    \left( \frac{\partial \mathbf{K}_0}{\partial Q_{\alpha\beta}} \right)_{ij} = \frac{\sigma}{2} Q^{-1}_{i\alpha} Q^{-1}_{j\beta}.
\end{equation}
This equilibrium form is especially useful because the explicit dependence of $\mathbf{K}_0$ on $\mathbf{Q}$ is simple, so the response functions can be written directly in terms of $\mathbf{Q}^{-1}$. Then, the response function of $\mathbf{K}_0$ due to the node parameter change is simplified as
\begin{equation}
    \chi_{ij(\alpha)} = -\frac{\sigma}{2} Q^{-1}_{i\alpha} Q^{-1}_{j\alpha},
    \label{chiKodr_eq}
\end{equation}
and the response functions for the change in a connection weight gives
\begin{equation}
    \chi_{ij(\alpha\beta)} = -\frac{\sigma}{2} Q^{-1}_{i\alpha} Q^{-1}_{j\alpha}+\frac{\sigma}{2} Q^{-1}_{i\alpha} Q^{-1}_{j\beta},
    \label{chiKodW_eq}
\end{equation}
Notice that a noise perturbation usually breaks detailed balance. For a system originally at equilibrium, its perturbed state does not satisfy detailed balance. The equilibrium response formula of noise perturbations should be derived from Eq. (\ref{partialKo_sigma}). These results provide a simple reference form for comparison with the nonequilibrium steady-state response.
